% Source PDF: source/普通高考/2012/2012北京理.pdf, page 1, question 4.
% Original PNG reference: img/2012/beijing_science/q4_fig1.png.
% Node -> directed-edge list: 开始→(k=0,S=1)→判定 k<3；
% 是→S=S·2^k→k=k+1→回到判定入口；否→输出 S→结束。
% The loop return arrow enters the main stem before the unique downward test arrow.
% Solid segments: every connector below; dashed segments: none.
% Labels: 是/否 are attached directly to their branch segments with natural anchors.
\begin{tikzpicture}[
  x=1cm,y=1cm,
  >=stealth,
  line width=0.55pt,
  line cap=round,line join=round,
  flowbox/.style={draw,anchor=center,align=center,minimum height=.66cm,
    inner xsep=4pt,inner ysep=2pt,
    execute at begin node={\setlength{\parindent}{0pt}\noindent}},
  io/.style={flowbox,shape=trapezium,trapezium left angle=78,trapezium right angle=102,
    trapezium stretches=false,minimum height=.78cm},
  terminal/.style={flowbox,rounded corners=7pt,minimum width=1.35cm,
    minimum height=.70cm},
  decision/.style={flowbox,diamond,aspect=2.8,minimum width=3.75cm,
    minimum height=1.05cm,inner sep=2pt},
  branchlabel/.style={inner sep=0pt,anchor=center,align=center,
    execute at begin node={\setlength{\parindent}{0pt}}},
]
  \node[terminal] (start) at (0,10.00) {开始};
  \node[flowbox,minimum width=2.90cm] (init) at (0,8.85) {$k=0,S=1$};
  \coordinate (pretest) at (0,8.00);
  \node[decision] (test) at (0,5.55) {$k<3$};
  \node[flowbox,minimum width=2.70cm] (multiply) at (3.05,6.45) {$S=S\cdot 2^k$};
  \node[flowbox,minimum width=2.35cm] (increment) at (3.05,7.65) {$k=k+1$};
  \node[io,minimum width=2.20cm] (output) at (0,3.45) {输出 $S$};
  \node[terminal] (finish) at (0,2.05) {结束};

  \draw[->] (start.south) -- (init.north);
  \draw (init.south) -- (pretest);
  \draw[->] (pretest) -- (test.north);

  % 是 branch: multiply, increment, then loop back to the test-entry stem.
  \draw[->] (test.east) -- node[below] {是} (3.05,5.55);
  \draw[->] (3.05,5.55) -- (3.05,6.12) -- (multiply.south);
  \draw[->] (multiply.north) -- (increment.south);
  \coordinate (loop-top) at (3.05,8.10);
  \coordinate (loop-tip) at (0.45,8.00);
  \draw[->] (increment.north) -- (loop-top) -- (loop-tip);
  \draw (loop-tip) -- (pretest);

  % 否 branch is independent of the loop and leads only to output.
  \draw[->] (test.south) -- node[right] {否} (output.north);
  \draw[->] (output.south) -- (finish.north);

  \path[use as bounding box] (-0.25,10.38) rectangle (4.00,1.98);
\end{tikzpicture}
